%% %%%%%%%%%%%%%%%%%%%%%%%%%%%%%%%%%%%%%%%%%%%%%%%% %%
%% Metadados do trabalho
%% AVISO: Todos esses dados serão automaticamente convertidos para caixa alta onde necessário
%% %%%%%%%%%%%%%%%%%%%%%%%%%%%%%%%%%%%%%%%%%%%%%%%% %%

%% Tipo de Trabalho
%% - Monografia
%% - Artigo
%% - Dissertação (Mestrado)
%% - Tese (Doutorado)
%% - Relatório técnico
\tipotrabalho{Artigo}


%% Título
\titulo{Modelo Preditivo para Avaliação de Desempenho Acadêmico}

%% Autor
\autor{João Victor dos Santos Mendes}


%% Nome do Curso (usado para a Capa do CD)
\nomedocurso{Licenciatura em Computação}

%% Local de publicação
\local{Teresina}

%% Preâmbulo do trabalho
\preambulo{Trabalho de Conclusão de Curso (Artigo Científico) apresentado como exigência parcial para obtenção do diploma do Curso de Licenciatura em Computação do Instituto Federal de Educação, Ciência e Tecnologia do Piauí, Campus Teresina Zona Sul.}

%% Orientador
%% "M\textsuperscript{e}." = Abreviação oficial para "Mestre"
\orientador{Prof. Dr. Kelson Carvalho Santos}

% Remova o comentário % se houver coorientador(a)
%\coorientador{Prof.ª Francisca Ocilma Mendes Monteiro}


%% Data do Trabalho
\data{2026}

%% Nome da Instituição (para a capa)
\instituicao{INSTITUTO FEDERAL DE EDUCAÇÃO, CIÊNCIA E TECNOLOGIA DO PIAUÍ
\\
CAMPUS TERESINA ZONA SUL
\\
CURSO LICENCIATURA EM COMPUTAÇÃO}

%% Primeiro membro da banca examinadora
\membroum{Prof. Ms. Francisco Assis Ricarte Neto}

%% Segundo membro da banca examinadora
\membrodois{Prof. Es. Fabrício Barroso de Carvalho}

%% Terceiro membro da banca examinadora
\membrotres{}

%% Data da apresentação do trabalho
%% Se não souber a data da apresentação, utilize \underline{\hspace{3.5cm}}
%% Isso cria um sublinhado de 3.5cm, onde você pode escrever a data depois!
%\dataapresentacao{02 de Outubro de 2023}
\dataapresentacao{\underline{\hspace{1.0cm}}/\underline{\hspace{1.0cm}}/\underline{\hspace{1.75cm}}}